% ============================================================
% Chapter 4: Proposed Method: Specification and Design
%   Section: The Fact-Checking Pipeline
%     Subsection: Claim Selection
% Included from chapters/04-proposed-method/06-fact-checking-pipeline/fact-checking-pipeline.tex
% Source: docs/arquitectura-tecnica.md 2.4.2; claim_selector.py.
% ============================================================

\subsection{Claim Selection}
\label{subsec:claim-selection}

Not every candidate claim of the enrichment stage is worth checking.
Two controls, both in the embedding space, reduce them to the few that
carry the article:

\begin{itemize}
    \item \textbf{Semantic deduplication}: when two candidates state the
    same fact in different words (similarity above 0.92), only one
    survives.
    \item \textbf{An anchor band}: the remaining candidates are ranked
    by how load-bearing they are for the article's thesis --- how
    central they are to its headline and lead, how much they share its
    main subjects, how specific they are --- and between two and four
    are kept. The ceiling stops an unusually dense article from
    multiplying the cost of verification; an article that yields fewer
    than the minimum is reported as such rather than padded.
\end{itemize}

Sentences that read as opinion or interpretation are dropped before the
ranking: no evidence can confirm or refute ``a hopeful step for the
sector''. The selector's precision --- the share of its proposals a
person accepts as checkable claims --- is measured by the daily
labelling batches, which ask the annotator to label or reject every
claim it proposes, with a reason (Section~\ref{sec:custom-dataset}).
